\section{Планковская решётка}
\label{sec:planck-lattice}

\begin{definition}[Планковская решётка]
\label{def:planck-lattice}
Носитель $\carrierLattice$ --- \latGCK\ $\mathbb{R}^3$
(теорема~\ref{thm:fcc-choice}, \cref{def:packing-types}).
Срез плоскостью $\{111\}$ даёт гексагональный слой $(2{+}1)$ (\cref{sec:hex-slice}, \cref{fig:carrier-fcc-111-slice}).
Объём элементарной ячейки Вороного
\[
  v_{\hv}=\caStepSpace^3/\sqrt{2}.
\]
\end{definition}

\begin{figure}[htbp]
  \centering
  \begin{tikzpicture}[carrier panel, iso view, scale=1.0]
  \def\r{0.28}
  \IsoSphereInk{0}{0}{0}{\r}
  \foreach \x/\y/\z in {
    1,1,0/1,-1,0/-1,1,0/-1,-1,0,
    1,0,1/1,0,-1/-1,0,1/-1,0,-1,
    0,1,1/0,1,-1/0,-1,1/0,-1,-1
  } {
    \IsoSphereMuted{\x}{\y}{\z}{\r}
  }
  \draw[carrier muted, dashed, fill=gray!8, opacity=0.45]
    (-1.2,-1.2,1.2) -- (1.2,-1.2,0) -- (1.2,1.2,-0.4) -- (-1.2,1.2,0.4) -- cycle;
  \node[carrier muted, font=\scriptsize, anchor=south west] at (-1.0,0.5,0.5) {$\{111\}$};
\end{tikzpicture}

  \caption{Носитель \latGCKshort\ и срез плоскостью $\{111\}$: вложенный гексагональный слой.}
  \label{fig:carrier-fcc-111-slice}
\end{figure}

Для этого носителя $\htick=\caKappaGeom\, t_P$ при $\caKappaGeom=1/\sqrt{2}$,
где $t_P=\caStepSpace/c$ --- планковское время в учебниковой нормировке; тогда $\caSignalSpeed=\sqrt{2}\,c$.
Это не подгонка задним числом: непрерывное определение $t_P=\caStepSpace/c$ предполагает $\caSignalSpeed=c$,
тогда как на решётке с $\caKappaGeom=1/\sqrt{2}$ тактовая и макроскопическая скорости
принципиально различаются.

